\section{An all-width consequence for linear activations}
\label{sec:linear}

The width qualification can be removed if the activation is changed to
$\phi(s)=s$. This is a separate result for a deep linear network; it is not
an extension to a wider tanh network. We retain the Boolean domain, the
full-cube output-margin restriction, exact label conditioning, and the
usual layer-scaled independent Gaussian prior.

Let all $L-1$ hidden layers have any common width $N\ge1$. Write
\[
 z_\theta(\xi)=vW^{L-1}\cdots W^1(1,\xi)
              =A_\theta\mathbin{\cdot}(1,\xi),
\]
where $W^1\in\mathbb R^{N\times(d+1)}$, the intermediate matrices are
$N\times N$, and $v\in\mathbb R^{1\times N}$. Independently, the entries
have variances
\[
 \operatorname{Var}(W^1_{ij})=\frac Nd\sigma^2,\qquad
 \operatorname{Var}(W^j_{i\ell})=\sigma^2\quad(2\le j<L),\qquad
 \operatorname{Var}(v_i)=\frac1N\sigma^2.
\]
The quadratic energy is
\[
 \kappa_N(\theta)=\frac dN\|W^1\|_F^2+
       \sum_{j=2}^{L-1}\|W^j\|_F^2+N\|v\|_2^2.
\]
For $N=1$ this agrees with the scalar prior used in the main theorem.

\begin{proposition}[Exponential mistakes at every sufficient width for
linear activations]\label{prop:linear-all-width}
For the target and fixed input permutation in
Theorem~\ref{thm:main}, the same $T_k=2^k-k-1$ forced rounds occur in the
strong-prior limit of this deep linear learner, at every fixed
$L\ge2$ and every fixed width $N\ge1$. All three conclusions of that
theorem hold at every sufficiently small positive scale; the scale
threshold may depend on $N$ as well as $k,L,\eta$.

At depth $L=k+2$, the target has a strict-margin linear-network
implementation with width one and coefficient bound two. Its
bounded-width-and-weight complexity for linear activation is exactly
$C_L^{\mathrm{lin}}(f)=2$. Hence the lower bound holds, in particular,
for every learner width $N\ge C_L^{\mathrm{lin}}(f)$, with no upper bound
on $N$.
\end{proposition}
\begin{proof}
For any desired effective affine vector $A$, the exact minimum energy
of its matrix factorizations is
\begin{equation}\label{eq:linear-factor-energy}
 \min_{\theta:A_\theta=A}\kappa_N(\theta)
       =L\bigl(d\|A\|_2^2\bigr)^{1/L}.
\end{equation}
To prove the lower bound, denote the $L$ factors by $M_1=W^1$,
$M_j=W^j$ for $2\le j<L$, and $M_L=v$, and put
$c_1=d/N$, $c_L=N$, and $c_j=1$ for the intermediate layers. Matrix norm
submultiplicativity and arithmetic--geometric mean give
\[
 \kappa_N=\sum_{j=1}^L c_j\|M_j\|_F^2
 \ge L\left(\prod_{j=1}^L c_j\|M_j\|_F^2\right)^{1/L}
 \ge L\bigl(d\|A\|_2^2\bigr)^{1/L}.
\]
Equality is attained by balanced rank-one factors. For $A\ne0$, choose
one unit vector $e\in\mathbb R^N$ and set
\[
 W^1=s_1e\frac{A}{\|A\|_2},\qquad
 W^j=s_j ee^\top\ (2\le j<L),\qquad v=s_Le^\top,
\]
where
\[
 r=(d\|A\|_2^2)^{1/(2L)},\qquad s_j=r/\sqrt{c_j}.
\]
Their product is $A$, since
$\prod_j s_j=r^L/\sqrt d=\|A\|_2$, and each energy summand is $r^2$.
For $A=0$, zero factors give equality as well.

The least squared norm of an affine vector with signed output margin
at least $1/2$ for a threshold function $g$ is $Q(g)/4$.
Equation~\eqref{eq:linear-factor-energy} therefore yields the exact
function energy infimum
\[
 E_{N,L}^{\mathrm{lin}}(g)
       =L\bigl(dQ(g)/4\bigr)^{1/L}.
\]
This expression is strictly increasing in $Q(g)$ and independent of
$N$. Strict and weak margins have the same infimum by slightly
increasing the final readout. The Gaussian small-noise argument from
Lemma~\ref{lem:energy} applies to these finite-dimensional matrix
parameters: the strict function cell is open, its closure is the
weak-margin cell, and the energy is coercive. Thus the posterior
concentrates on the unique minimum-$Q$ consistent function after each
specified history. Lemma~\ref{lem:minimum} and the same finite-round
union bound prove the claimed mistake statements at every fixed width.
No interchange of the width and small-scale limits is needed.

For the teacher, divide $A_{2^k-1}$ by $2^k$ in the first layer, use
$k$ further hidden scalar layers of weight two, and use readout weight
one. There are $k+1$ hidden layers, so $L=k+2$, and the final output is
exactly $A_{2^k-1}$, with absolute margin at least one. All weights are
bounded by two. Zero-padding to width two proves
$C_L^{\mathrm{lin}}(f)\le2$.

If this complexity were one, the product of the scalar downstream
weights would have absolute value at most one. A strict output margin
above $1/2$ would then force a first affine margin $\rho>1/2$.
Normalize that affine vector by $\rho$ and apply the diagonal and carry
inequalities of Lemma~\ref{lem:minimum}. They force its original
$w_{k-1}$ to exceed $2^{k-2}>1$, contradicting the coefficient bound
one. Hence $C_L^{\mathrm{lin}}(f)=2$.
\end{proof}

This linear-activation result gives a polynomial-iid versus
exponential-online separation even at arbitrarily large learner widths:
the represented Boolean functions are still exactly the affine
threshold functions, so the iid argument in the preceding section is
unchanged. For tanh, additional neurons introduce nonlinear classifiers,
and the factorization identity above no longer applies. The result
with tanh and $N\ge C_L(f)$ therefore remains open.
